The ontology of \autoref{sec:appendixA} closes the loop from the signal to the agent decision by making the tokenisation geometry a first-class object rather than a property of the model: a \pv{TimeSeriesSignal} carries one \pv{PatchedRepresentation} per geometry $(P,S)$ under evaluation, each recording its patch size, stride and overlap. A frequency-specific loss of observability therefore attaches to the signal-configuration pair rather than to Chronos-Bolt-Tiny alone, so the same checkpoint can be safe for one plant and not for another. Because the lock conditions of \autoref{eq:locking-candidate} are arithmetic on those same fields, the candidate set is derivable from the knowledge base rather than hard-coded beside it: a reasoner can list the sites of a geometry never measured, and recompute them whenever the geometry or the sampling interval changes, a checkpoint change mattering only when it alters $(P,S)$.

What that arithmetic reaches is the domain vocabulary, none of whose terms is read off the raw telemetry. A \valueof{RampEvent}, or a \valueof{PartialCloud} generation context, is inferred from an observation whose short-horizon forecast the agent has already trusted, so a frequency the representation cannot hold is one whose signature cannot support that inference. The exposure is selective and the arithmetic says where: on one-minute telemetry with $P=16$ the patch comb falls at periods of $16/k$ minutes, the timescales of cloud-edge crossings and inverter duty cycling, while a broadband \valueof{InverterTrip} is untouched, since no comb hides a step. The predicted failure is therefore a missed distinction rather than a wrong number: two contexts differing only at a candidate frequency collapse to one latent state.

For monitoring, placement matters more than cadence. The stride-derived and patch-derived candidates follow separately from $\mathcal{F}_{S}$ and $\mathcal{F}_{P}$, but neither can be watched downstream: a detector placed after tokenisation inherits the blind spot it is meant to find, so the monitor acts on the raw telemetry, the one stage the forecaster cannot perform itself.

The explanation the agent produces keeps a geometric candidate apart from an established loss. Each \pv{ObservabilityAssessment} records the estimand, the effect size, the interval and the posterior probability behind its verdict, and assigns \valueof{StructurallyAliased} only to a frequency that is a lock candidate and also carries confirmed evidence of loss; \autoref{eq:bayesian-evidence} separates the two directions of evidence, and because $\tau>1/2$ they cannot both hold, so everything short of confirmed preservation is \valueof{PartiallyObservable}. One line per assessment reaches a person: frequency, geometry, estimand with interval, state, and the control authority that follows. Its arithmetic half is derivable for any frequency, its measured half only where a fit exists, and a missing fit is reported rather than defaulted away.

Action enters control through \autoref{eq:mode-selection} and not through a policy of its own, so the assessment rather than the geometry selects the mode. A confirmed \valueof{StructurallyAliased} state selects \valueof{SafeMode} at the same priority as a sensor malfunction, under \autoref{eq:priority}, the array staying connected because the fault is in the model and not in the plant; a merely \valueof{PartiallyObservable} state selects \valueof{ConservativeMode}, which lowers authority without suspending dispatch. Both are reversible, degrading authority rather than substituting a correction the analysis does not support, and the default runs in the safe direction: where no fit covers a candidate the state is \valueof{PartiallyObservable} by construction, so an unmeasured plant runs permanently at reduced authority and pays the cost in capacity rather than in risk.
